\documentclass[a4paper]{article}

\usepackage{array}
\usepackage{fullpage}
\usepackage{enumerate}
\usepackage{hyperref}
\usepackage[all]{hypcap}
\usepackage{listings}
\usepackage{color}
\usepackage{graphicx}

\definecolor{light-gray}{gray}{0.95}
\definecolor{dark-green}{rgb}{0, 0.5, 0}
\definecolor{dark-gray}{gray}{0.4}
\definecolor{Gray}{gray}{0.95}

\lstset{ %
	language = C,                   % choose the language of the code
	basicstyle = \small\ttfamily,   % the size and fonts that are used
	frame = single,                 % adds a frame around the code
	tabsize = 3,                    % sets default tabsize
	breaklines = true,              % sets automatic line breaking
	numbers = left,                 % where to put the line-numbers
	numberstyle = \footnotesize,    % the style of the line-numbers
	backgroundcolor = \color{Gray}, % the background color of the listing
	showstringspaces=false,
	keywordstyle=\color{blue},
}

\newcolumntype{V}{>{\centering\arraybackslash} m{.75\linewidth} }

\title{Assignment 3: MPI}
\author{Jacco Spoelder (s1348493) \and Xeryus Stokkel (s2332795)}

\begin{document}

\maketitle

\section{Exercise 1}
\begin{enumerate}[(a)]
	\item When cutting the interval into equal chunks we obtain the results in \autoref{tbl:prime}. The code to share the work can be found in \autoref{code:prime1sharing}. The master process splits the entire $[a, b]$ interval into equally sized sections and sends them out to all the different slaves. After this each process will work on its sub-interval, we didn't have to change any code to do this since the original code also worked with calculating ranges. Finally the results are shared via a simple \texttt{MPI\_Reduce} and the master thread outputs the result as \autoref{code:prime1finish} shows.
	\lstinputlisting[firstline=33,lastline=59, label={code:prime1sharing},caption={Work sharing code for the static domain decomposition}]{prime/prime.c}
	\lstinputlisting[firstline=80,lastline=82, label={code:prime1finish}, caption={Result sharing code for the static domain decomposition}]{prime/prime.c}
	
	From \autoref{tbl:prime} we can see that when using multiple processors we can see that there is a discrepancy in the runtime for different processes. The longest runtime can be four times as high as the lowest runtime. The cause of this is quite obvious, if there are two processes then one will work on the interval $[1, \ldots, 2.5 \cdot 10^7]$ while the other will work on the interval $[2.5 \cdot 10^7, \ldots, 5 \cdot 10^7]$. This means that the first process will need to check roughly $\sqrt{\sum_{n=1}^{2.5 \cdot 10^7}n} \approx 1.7 \cdot 10^7$ numbers, the second process needs to check $\sqrt{\sum_{n=2.5\cdot 10^7}^{5\cdot 10^7}n} \approx 3 \cdot 10^7$ numbers in total in the \texttt{isPrime} function. This is almost a doubling of the amount of possible divisors to check and we can see that the runtime is also almost double from \autoref{tbl:prime}. Once the process with the largest numbers in its interval is done the program will also complete after displaying the total number of primes.
	
	\item We used the master-worker model as our parallelization strategy. The master process sends out small parts of the interval to different workers which will then work on calculating the number of primes. When all work has been distributed a termination signal will be send to each worker and the worker processes will send their results back to the master process upon receiving the termination signal. The code can be found in \autoref{code:prime2}.
	
	Special care has been taken to ensure the ordering of messages to ensure that workers don't accidentally ignore work because they received the termination signal early. This happened during some of the early test runs so we took the time to fix this without having to fundamentally change how our program works.
	
	Because of the way that communication works the master process doesn't actually do any work on finding the number of primes. Therefore there is no data for 1 process in \autoref{tbl:prime2} since it would never complete. From \autoref{tbl:prime2} we can see that the speed-up is a lot higher than in \autoref{tbl:prime}. It should be noted that there is effectively one process less working on the problem than the table states so the achieved speed-up would be higher if all the processes would've been working on the problem. In \autoref{tbl:prime2} we also see that the difference in highest and lowest runtime is about 100ms so the load is a lot better balanced amongst the processors as none of them finishes while another process has a lot more work to do.
\end{enumerate}

\begin{table}[h]
	\centering
	\caption{Runtimes for finding the number of primes by splitting the load into equal sized parts. Runtime for the sequential program is 2:22.72, the speed-up is based on comparing to this time.}
	\label{tbl:prime}
	\begin{tabular}{l|r|r|r|r}
		Processes & Lowest time & Highest time & Total time & Speed-up \\ \hline
		 1 &  & 2:19.60 & 2:19.74 & 1.02 \\
		 2 & 50.61 & 1:28.64 & 1:28.78 & 1.61 \\
		 4 & 18.43 & 47.92 & 48.20 & 2.96 \\
		 8 &  6.74 & 24.71 & 26.16 & 5.46 \\
		12 &  3.74 & 16.67 & 19.36 & 7.37 \\
	\end{tabular}
\end{table}

\begin{table}[h]
	\centering
	\caption{Runtimes for finding the number of primes by interleaving the load. Runtime for the sequential program is 2:22.72, the speed-up is based on comparing to this time.}
	\label{tbl:prime2}
	\begin{tabular}{l|r|r|r|r}
		Processes & Lowest time & Highest time & Total time & Speedup \\ \hline
		 2 &    & 2:15.49 & 2:16.68 & 1.04 \\
	 	 4 & 45.53 & 45.62 & 46.85 & 3.04 \\
	 	 8 & 19.49 & 19.59 & 21.06 & 6.78 \\
		12 & 12.41 & 12.55 & 14.26 & 10.01 \\
	\end{tabular}
\end{table}

\section{Contrast stretching}
\begin{enumerate} [(a)]
	\item We created a function $masterProcess$ which is only called by the root process, to read the image file and distribute the workload to all the processes, including the root process itself. We use the MPI function $scatterv$ 
	for this. Scatterv sends only a part of the whole datablock,defined by the user-specified size, to all other processes instead of all data, which makes it more memory efficient. First however we need to determine the data to be send, and how large the blocks are. With the routine:
	\lstinputlisting[firstline=183,lastline=193,label={code:contrast},caption={Work division and sending}]{contrast/contrast.c}
	we divide the array with the pixeldata in equal chunks and send the size of those chunks to the slave processes. When the dividing does not yield equal chunks, the last chunk is shortened.
 
	With a $MPI_Allreduce$ operation:
	\lstinputlisting[firstline=199,lastline=203,label={code:contrast},caption={Calculate and broadcast minimal and maximum pixel values}]{contrast/contrast.c}
	We reduce the minimum and maximum pixel values and rebroadcast them to all processes, so each process holds its own datablock, but with the proper pixel values. Each process can now recalculate the pixelvalues of its own datablock.
	A function $slaveProcess/$ is created where all processes other than the root handles its data and calculations. Because the root scatters the data and reduces en rebroadcasts the minimum and maximum pixelvalues, for each slave process the sending and retreiving of this data has to be aranged, using the same functions $MPI_scatterv$ and $MPI_gatherv$. After contrast stretching, we can now retrieve the data to the root with a $MPI_Gatherv$ operation:
	\lstinputlisting[firstline=232,lastline=233]{contrast/contrast.c} 
	In the $slaveProcess$, the function $MPI_Gatherv$, is also used to send the results to the root. The root is now able to write the image.
	 \item The MPI function $MPI_Wtime()$ can be used to calculate the times each process runs. It returns the elapsed time from some fixed point. Each process wil have its own instance. We measure the time after each process is finished calculating its data and has send it back tot the root: 
	\lstinputlisting[firstline=250,lastline=256]{contrast/contrast.c} 
	
	\item Runnig the contrast stretch program on the 12-core millipede cluster with the kdk.pgm image file gives us the following calculation times with different threads, for the smallest time a thread is working and the longest time a thread is working:

\begin{table}[h]
	\centering
	\caption{Runtimes for contrast stretching a pgm image.}
	\label{tbl:contrast}
	\begin{tabular}{l|r|r|r}
		Processes & Lowest time & Highest time(root) & Speed-up \\ \hline
		 1 & 12.3 & 12.3 & none       \\
		 2 & 5.6 & 10.2. & 20 percent \\
		 4 & 4.2 & 9.2   & 33 percent \\
		 8 &  3.5 & 8.5  & 44 percent \\
		12 &  3.2 & 8.3 &  48 percent\\
	\end{tabular}
\end{table}
The highest timings for a thread is always for the root process, since this is the one starting, dividing, integrating the workload and outputting the end result. Increasing the number of processes by far does not linearly scale with the corresponding maximum time.The cause is that only a part of the work is parallelized. The root process also takes advantage of this, but the difference is quite small, especially with larger number of threads. Also the communication to and from processes take more time with increasing numbers of processes.The main workload which is parallelized(lowest times) scale better with increasing processes, since this part is almost fully parallelized. Communication oerhead also impedes the speed here however. 

\end{enumerate}
\section{Mandelbrot fractal}
\begin{enumerate}[(a)]
	\item Because the exact size of the output is known we can simply calculate which part a process has to do based on its rank. We decided to split the image up into horizontal sections since calculating the image is done in a scanline fashion. Each thread will process a block of lines in the image. Determining which part of the image a process needs to calculate is shown in \autoref{code:mandelbrot1split}. We also needed to adjust the \texttt{mandelbrotSet} function since it was not designed to be ran in parallel, if we wouldn't do this then each process would calculate the entire image instead of just a small section. The code for this function can be found in \autoref{code:mandelbrot1set}.
	\lstinputlisting[firstline=150,lastline=157,label={code:mandelbrot1split},caption={Determining the section of the Mandelbrot image a thread needs to calculate.}]{mandelbrot/mandelbrot.c}
	\lstinputlisting[firstline=109,lastline=136,label={code:mandelbrot1set},caption={The \texttt{mandelbrotSet} function}]{mandelbrot/mandelbrot.c}
	
	Finally we need to merge all the partial images. This is done line by line since we didn't adjust the way images are stored to send blocks of the data. The code to merge is found in \autoref{code:mandelbrot1merge}
	\lstinputlisting[firstline=159,lastline=171,label={code:mandelbrot1merge},caption={Merging the results of the Mandelbrot Set calculating processes.}]{mandelbrot/mandelbrot.c}
		
	Runtimes are displayed per process in \autoref{tbl:mandelbrot}. Each column shows the runtime for each process for a given number of processors. The total runtime for the program is displayed at the bottom. We can see that the columns follow a certain symmetry: the first and last process have the lowest running time while the processes in the middle have the highest running time. The difference in running time is very large when the number of processes grows!
	
	This phenomenon is caused by how we distribute the work. We split the image up in $N$ rows (where $N$ is the number of processes) that are of equal size. From \autoref{fig:mandelbrot} we can see that the complexity at the top and the bottom of the fractal is a lot lower than the complexity in the middle. \autoref{tbl:mandelbrotcalc} shows the image split up into 8 sections, listing the number of iterations required to calculate each section. From this we can see that first and last process have a lot less work to do ($x^2 + y^2 > 10$ is satisfied in fewer iterations). A better solution would be to distribute work in such a way that each process has to deal with the same amount of complexity.
	
	\item To better balance the load across the different processes we opted for the master-slave model. Process 0 is now a process that only works on distributing work to the various slaves. This means that there must be at least 2 processes to run the program, otherwise the master will endlessly wait until a slave requests some work but this will never happen because there are no slaves if there is just 1 process.
	
	We use a very simple model, a slave returns the data for a line in the image, the master stores it and sends the number of a new line to work on. This process repeats until all lines have been send out, in this case the master sends a kill signal to the slaves so they know they can stop. When all slaves have been send a kill signal the master will also stop. This leaves just one problem, the slaves don't know which line they should start to work on since they will only get a line number in response to delivering a result. This was solved by having the slaves send an invalid line number, which the master process will not store in the final image but it will send a command to the slave process so it knows on which line to work. The code that does this can be found below in \autoref{code:mandelbrot2sharing}. We also removed the outer for-loop from the \texttt{mandelbrotSet} function since it is not needed anymore as it will operate on a line-by-line basis. It now takes an extra argument that specifies which line should be calculated.
	
	\lstinputlisting[firstline=193,lastline=236,label={code:mandelbrot2sharing},caption={Load balancing code for the Mandelbrot Set calculation}]{mandelbrot/mandelbrot2.c}
	\lstinputlisting[firstline=153,lastline=177,label={code:mandelbrot2set},caption={The function which calculates a part of the image has also been changed.}]{mandelbrot/mandelbrot2.c}
	
	The runtime results for this code are found in \autoref{tbl:mandelbrot2}. Just like with \autoref*{tbl:prime2} there is no time for 1 process because the master doesn't do any work itself. As we can see from \autoref{tbl:mandelbrot2} the times for each process are a lot better balanced. Most processes take an equal amount of time to complete, with the exception of the master process. This is because the time to write the image is also taken included in this time so it will always take longer than the threads that only need to do the calculations.
\end{enumerate}

\begin{table}[h]
	\centering
	\caption{Runtimes for generating an Mandelbrot set image using static domain decomposition. Process times are the time that each process takes to calculate its part of the image, while total time is the time that the program took including memory allocation, MPI communication and writing the fractal to a file.}
	\label{tbl:mandelbrot}
	\begin{tabular}{l|r|r|r|r|r}
		Number of processes & 1 & 2 & 4 & 8 & 12 \\ \hline
		process 0  & 40.76 & 20.36 &  1.14 &  0.06 &  0.03 \\
		process 1  &       & 20.40 & 19.19 &  1.10 &  0.12 \\
		process 2  &       &       & 19.25 &  7.18 &  1.04 \\
		process 3  &       &       &  1.17 & 12.03 &  4.36 \\
		process 4  &       &       &       & 12.06 &  6.20 \\
		process 5  &       &       &       &  7.20 &  8.70 \\
		process 6  &       &       &       &  1.12 &  8.71 \\
		process 7  &       &       &       &  0.06 &  6.18 \\
		process 8  &       &       &       &       &  4.37 \\
		process 9  &       &       &       &       &  1.03 \\
		process 10 &       &       &       &       &  0.11 \\
		process 11 &       &       &       &       &  0.04 \\ \hline
		Total      & 41.63 & 21.03 & 20.92 & 13.95 & 10.19
	\end{tabular}
\end{table}

\begin{figure}[h!]
	\centering
	\includegraphics[width=.75\textwidth]{mandelbrot/mandelbrot.png}
	\caption{Final render of the Mandelbrot fractal}
	\label{fig:mandelbrot}
\end{figure}

\begin{table}[h]
	\centering
	\caption{8 sections of the Mandelbrot fractal including and the number of iterations required to compute each section. The intensity of a pixel represents the amount of iterations needed to compute its value. A black pixel means no calculations while a white pixel means that the iteration limit has been reached.}
	\label{tbl:mandelbrotcalc}
	\begin{tabular}{Vr}
		\includegraphics[width=.75\textwidth]{mandelbrot/mandelbrotcalc0.png} & 6 123 840 \\
		\includegraphics[width=.75\textwidth]{mandelbrot/mandelbrotcalc1.png} & 183 359 327 \\
		\includegraphics[width=.75\textwidth]{mandelbrot/mandelbrotcalc2.png} & 1 218 769 786 \\
		\includegraphics[width=.75\textwidth]{mandelbrot/mandelbrotcalc3.png} & 2 042 471 058 \\
		\includegraphics[width=.75\textwidth]{mandelbrot/mandelbrotcalc4.png} & 2 046 873 095 \\
		\includegraphics[width=.75\textwidth]{mandelbrot/mandelbrotcalc5.png} & 1 221 147 602 \\
		\includegraphics[width=.75\textwidth]{mandelbrot/mandelbrotcalc6.png} & 184 863 811 \\
		\includegraphics[width=.75\textwidth]{mandelbrot/mandelbrotcalc7.png} & 6 129 020 \\
	\end{tabular}
\end{table}

\begin{table}[h]
	\centering
	\caption{Runtimes for generating an Mandelbrot set image by using the master-slave model. Process time is the time that a process took for MPI communication and fractal calculation while total time includes setup and writing the obtained image to a file.}
	\label{tbl:mandelbrot2}
	\begin{tabular}{l|r|r|r|r}
		 & 2 & 4 & 8 & 12 \\ \hline
		process 0  & 41.26 & 14.05 & 6.30 & 4.16 \\
		process 1  & 40.83 & 13.62 & 5.84 & 3.73 \\
		process 2  &       & 13.62 & 5.84 & 3.73 \\
		process 3  &       & 13.62 & 5.84 & 3.73 \\
		process 4  &       &       & 5.84 & 3.74 \\
		process 5  &       &       & 5.84 & 3.72 \\
		process 6  &       &       & 5.84 & 3.73 \\
		process 7  &       &       & 5.84 & 3.73 \\
		process 8  &       &       &      & 3.73 \\
		process 9  &       &       &      & 3.73 \\
		process 10 &       &       &      & 3.73 \\
		process 11 &       &       &      & 3.73 \\ \hline
		Total      & 42.44 & 15.30 & 7.71 & 5.92
	\end{tabular}
\end{table}

\section{Volume rendering}
\begin{enumerate}[(a)]
	\item To make a parallel version of the program we had to do several things. The first was changing the way that images and 3D data is stored in memory so that it is easier for MPI to send it between different processes. The new structures are shown in \autoref{code:renderstruct}. The code for allocating and freeing memory is shown in \autoref{code:rendermalloc}. Most of this code can also be found on the hints page on Nestor, but we had to roll our own deallocation routines.
	
	\lstinputlisting[firstline=25,lastline=39,label={code:renderstruct},caption={Updated structures to support sending data between processes}]{volumerender/render.c}
	\lstinputlisting[firstline=59,lastline=108,label={code:rendermalloc},caption={Updated memory allocation and deallocation routines to work with the new memory structure.}]{volumerender/render.c}
	
	Rendering a frame of the animation doesn't depend on the last frame but only on the number of the frame and the volume data. Because of this it is easiest to parallelize the program by making threads rendering single frames. Again we went for a master-slave model where the master does nothing but delegate the work to the slaves, this means that at least two processes are necessary to complete the task.
	
	The master starts by allocating memory to keep track of which slave is working on which frame. Results are received asynchronously so there is also a need to keep track of whether a slave has send its result. After this is done the master process reads the volume data and sends it to the slave processes. The code for allocating and sending the data is shown in \autoref{code:rendermasterinit}, the slave process code for receiving the dat is shown in \autoref{code:renderslaveinit}.
	
	\lstinputlisting[firstline=351,lastline=372,label={code:rendermasterinit},caption={Master process code for allocating memory for bookkeeping and loading and distributing the volume data}]{volumerender/render.c}
	\lstinputlisting[firstline=417,lastline=433,label={code:renderslaveinit},caption={Slave process code for retrieving the volume data.}]{volumerender/render.c}
	
	Next the master needs to initialise the asynchronous communication and at the same time it sends out the initial frames for the slave processes to work on (see \autoref{code:rendermasterrecv}). Note that we had to start receiving on the \texttt{MPI\_Request} belonging to the master because \texttt{MPI\_Waitany} doesn't support uninitialised requests. We then have to make sure that no initial frame is send to the master process because it doesn't do any work (and it would also cause the master process to freeze).
	
	\lstinputlisting[firstline=374,lastline=383,label={code:rendermasterrecv},caption={The master process starts receiving data asynchronously and hands out the first tasks.}]{volumerender/render.c}
	
	The master then starts its main loop where it waits for a slave process to return a result. The first thing that is done is responding to the result by handing a new task or by sending a kill signal to the slave process when all the work has been handed out. After doing this the master process writes the received data to file. At first this was done the other way around, but this led to a small speed penalty because a slave process would have to wait for the master process to finish writing the file before it could start working on the next frame. Now a slave can start doing more work right away. The code that orchestrates all of this can be viewed in \autoref{code:rendermastermainloop}.
	
	\lstinputlisting[firstline=385,lastline=407,label={code:rendermastermainloop}, caption={The master handles result data and sends out a new task at the same time.}]{volumerender/render.c}
	
	In the mean time the slave processes calculate frames as long as tasks are being received. They do no asynchronous communication because they only communicate with the master process so there is no harm in having to wait for communication or work to finish since they can't do anything else in the mean time. The main loop for slave processes is shown in \autoref{code:renderslavemainloop}
	
	\lstinputlisting[firstline=436,lastline=453,label={code:renderslavemainloop}, caption={The slave calculates frame data as tasks are being received.}]{volumerender/render.c}
	
	After this only clean-up is needed for all processes.
	
	\item Timing results for the $128 \times 128 \times 128$ input file can be found in \autoref{tbl:render}. The first thing of note is that there is a small speed-up with just two threads, this is slightly odd because there is still only one process doing all the calculations. The master process is doing all the disk IO though, so the thread which is calculating the data doesn't have to wait for a hard drive before it can start calculating again. This small concurrency is enough to be a small enough speed-up. The speed up between 8 and 12 processes is not very large, this is most likely due to processes needing to wait for the master to write another processes' result to the disk before it will hand out a new task to that process. This means that slaves are waiting on the master thread a sizeable portion of the time.
	
	Another remarkable thing is that the change in speed up going from 4 to 8 processes is very large! There doesn't seem to be a good explanation for this. The most likely cause is that with 8 processes the communication between the master and the slaves is just rightly balanced so that slaves don't really have to wait for each other to finish communicating with the master, they also don't have to wait that long for the master to finish writing the received data to a file.
	
	\item When the input data is $256 \times 256 \times 256$ voxels large the runtimes in \autoref{tbl:render2} are obtained. What is most notable in comparison to \autoref{tbl:render} is that the speed up is lower is that the large jump in speed up between 4 and 8 processes is mostly gone. This is most likely because the amount of time to calculate a frame has increased more than the time required to write a frame to disk. This means that slave processes haven't hit a sweet spot between doing calculation and communicating/waiting for the master.
	
	The maximum speed up also seems to be higher when comparing the two differently sized inputs. This seems to be a case of Gustafson's law in working, when you increase the size of the dataset then you will also see an increase in speed up.
\end{enumerate}

\begin{table}[h]
	\centering
	\caption{Volume renderer run times for different number of processes. Process time is the maximum of the time that all subprocesses take to do their calculations. Total time is the time that the program ran, this includes setup and system time.}
	\label{tbl:render}
	\begin{tabular}{r|r|r|r}
		Processes & Process time & Total time & Speed up\\ \hline
		Sequential & & 48.05 \\
		 2 & 46.08 & 46.30 & 1.04 \\
		 4 & 21.05 & 22.27 & 2.15 \\
		 8 &  5.45 &  6.85 & 7.01 \\
		12 &  3.72 &  6.41 & 7.49
	\end{tabular}
\end{table}

\begin{table}[h]
	\centering
	\caption{Run time for the volume renderer when the input file is $256 \times 256 \times 256$ voxels large. Process time is the maximum of the time that all subprocesses take to do their calculations. Total time is the time that the program ran, this includes setup and system time.}
	\label{tbl:render2}
	\begin{tabular}{r|r|r|r}
		Processes & Process time & Total time & Speed up \\ \hline
		Sequential & & 5:39.97 & \\
		 2 & 5:30.26 & 5:30.41 & 1.03 \\
		 4 & 2:03.99 & 2:05.21 & 2.72 \\
		 8 &   52.60 &   54.02 & 6.29 \\
		12 &   37.75 &   39.48 & 8.61
	\end{tabular}
\end{table}

\section{Solving the wave equation}
As in some of the other programs of these assignments, we create a master process and slave process(es), where the master process distributes the work to the slave processes and also gathers the finished work from these slave processes. The main question here is which part, which data, is most efficient to parallelize. Since a standard or user defines number of frames have to be calculated, and each frame is dependent on the previous frame, we cannot divide a subset of the frames between the processes. Therefore we divide the calculation work for 1 frame between the processes. Since the size of the frame is determined by the number of grid cells $n$, we base the division of work on the grid size and distribute equal parts of it to the processes. The other settings for the frames are the same for each process, so we broadcast those to all processes with the following code:
\lstinputlisting[firstline=359,lastline=362,label={code:wavesettings},caption={Send all settings to slaves}]{wave/wave.c}
The worksize is determined based on $n$ and divided equally. If the worksize cannot be split in equal chunks, the last process gets the smallest part. Buffers of the correct size for receiving the work are also allocated:
\lstinputlisting[firstline=364,lastline=380,label={code:wavesettings},caption={Work division and sending}]{wave/wave.c}
Then each process has to calculate its part by calling the $solveWAve$ and $contrastStretch$ function. The master process does this in the $masterProcess$ function and the slaves in the $slaveProcess$ function  Also a timer is used to obtain the time each process is takes to solve the wave equation. 
\lstinputlisting[firstline=386,lastline=396,label={code:wavesettings},caption={Call for wave solvation and contrast stretching}]{wave/wave.c}
Each process which has solved its part of the wave equation and the contrast stretching of a particular frame, copies it into an allocated block of memory called $data$ at the right position, after which al the data of one frame is send back to the all processes with $MPI_Gatherv$, so each process is able to calculate its part for the next frame, based on the previous one. For the solved wave:
\lstinputlisting[firstline=167,lastline=169,label={code:wavesettings},caption={Send solved wave data to other processes}]{wave/wave.c}
And for the contrast stretching:
\lstinputlisting[firstline=205,lastline=207,label={code:wavesettings},caption={Send contrast stretch data to other processes}]{wave/wave.c}
After all frames are calculated, the master process saves the frames with the $saveFrame$ function. All the used memory is released then.


Timing results, averages of 5 runs on the 12-core millipede cluster, can be found in \autoref{tbl:wavetime}

\begin{table}[h]
	\centering
	\caption{Runtime for solving the wave equation, with average time per process and total duration including saving all frames by the master process.}
	\label{tbl:wavetime}
	\begin{tabular}{r|r|r|r}
		Processes & Process time & Total time & Speed up \\ \hline
		 1 & & 9.2 & \\
		 2 & 2.0 & 7.0 & 31\% \\
		 4 & 1.8 & 8.5 & 8\% \\
		 8 & 1.5 & 8.0 & 15\% \\
		12 & 1.1 & 6.0 & 53\% 
	\end{tabular}
\end{table}

The speedup on the millipede cluster is somewhat irregular for the runs with multiple processes. The difference in timing for each process, doing the calculations, are quite scalable. With each increase of processes, the calculation time of each separate non-root process becomes smaller. The total timing, where the root process is also saving the frames is more irregular, but times are smaller with larger numbers of processes. We also run this on a home computer with 4-cores(data not shown), which shows the same pattern with the time for each process, however the total time only decreases from 1 to 2 processes. With 3 and 4 processes total times is equal as when run with 2 processes. These effects are probably seen because of the increase in the amount data which has to distributes as the number of processes increase. We also noted a increasing memory usage with increasing number of processes.





\clearpage
\appendix
\section{Full program code}
\subsection{Program code for prime number search}
\lstinputlisting[label={code:prime1}, caption={Code for prime search by static domain decomposition}]{prime/prime.c}
\lstinputlisting[label={code:prime2}, caption={Code for prime search using load balancing}]{prime/prime2.c}

\subsection{Code for contrast stretching}
\lstinputlisting[label={code:contrast}, caption={Code for contrast stretching}]{contrast/contrast.c}

\subsection{Code for mandelbrot fractal}
\lstinputlisting[label={code:mandelbrot}, caption={Code for calculating the Mandelbrot fractal using static domain decomposition}]{mandelbrot/mandelbrot.c}
\lstinputlisting[label={code:mandelbrot2}, caption={Code for calculating the Mandelbrot fractal using one process to hand out work and the others doing the actual work}]{mandelbrot/mandelbrot2.c}

\subsection{Code for volume rendering}
\lstinputlisting[label={code:volume}, caption={Code to render the rotation of a volume}]{volumerender/render.c}

\subsection{Code for solving the wave equation}
\lstinputlisting[label={code:wave}, caption={Code to solve the wave equation in parallel}]{wave/wave.c}

\end{document}